% =============================================================================
%                    WEEK 4: REASONING AND COMMUNICATION
% =============================================================================
\section{Week 4: Reasoning and Communication}

% -----------------------------------------------------------------------------
%                              4.1 TOPIC SUMMARY
% -----------------------------------------------------------------------------
\subsection{Topic Summary}

\begin{keybox}[Learning Objectives]
By the end of this week, you will be able to:
\begin{enumerate}
  \item Explain \textbf{logic-based (symbolic) AI} and its foundations.
  \item Describe \textbf{machine learning} approaches and their successes.
  \item Identify \textbf{limitations} of both paradigms.
  \item Understand \textbf{non-monotonic reasoning} and why it matters.
  \item Explain \textbf{argumentation theory} as a formalization of reasoning.
  \item Describe how \textbf{dialogue} enables distributed reasoning among multiple agents.
\end{enumerate}
\end{keybox}

\separator

\subsubsection*{The Big Picture}

AI has two major paradigms: \textbf{logic-based (symbolic)} and \textbf{machine learning}. Each has strengths and limitations. Logic-based AI struggles with real-world complexity and brittleness; ML struggles with explainability, causation, and common-sense reasoning. \textbf{Argumentation theory} offers a middle ground---formalizing non-monotonic reasoning in a way that is understandable to humans, computationally tractable, and supports \textbf{joint reasoning via dialogue}. This has important ethical implications: transparent, collaborative AI-human reasoning.

\separator

\subsubsection*{What Examiners Like to Ask}

\begin{itemize}
  \item Compare logic-based AI and machine learning approaches.
  \item What is monotonic vs.\ non-monotonic reasoning? Why does it matter?
  \item Explain the symbol grounding problem.
  \item What are the limitations of machine learning (black box, no causation, etc.)?
  \item What is argumentation? How do arguments attack each other?
  \item Explain Dung's argumentation framework and the labeling approach.
  \item What are grounded and preferred extensions?
  \item How can argumentation support distributed reasoning via dialogue?
\end{itemize}

\bigseparator

% -----------------------------------------------------------------------------
%                 4.2 OVERVIEW + CORE CONCEPTS + EXTENDED THEORY
% -----------------------------------------------------------------------------
\subsection{Overview + Core Concepts + Extended Theory}

\subsubsection*{Roadmap}
\begin{enumerate}
  \item Logic-based (symbolic) AI
  \item Machine learning AI
  \item Limitations of both paradigms
  \item Non-monotonic reasoning
  \item Argumentation theory (Dung's framework)
  \item Distributed reasoning via dialogue
\end{enumerate}

\bigseparator

% --- Logic-based AI ---
\subsubsection{Logic-Based (Symbolic) AI}

\begin{defbox}[Good Old-Fashioned AI (GOFAI)]
Early AI used \textbf{symbolic formulae} to represent beliefs, desires, and goals. Reasoning involved applying \textbf{inference rules} to derive new facts about the world and determine appropriate actions.

\textbf{Key idea:} Describe the world in a logical language (like natural language describes states of affairs), then apply inference rules to deduce new truths and actions.
\end{defbox}

\paragraph*{Logic in Mathematics (Background)}
\begin{itemize}
  \item Start with \textbf{self-evidently true axioms} (e.g., $\alpha \lor \neg\alpha$).
  \item Apply \textbf{inference rules} (e.g., disjunctive syllogism: from $\alpha \lor \beta$ and $\neg\alpha$, infer $\beta$).
  \item Derive complex truths from simple axioms.
  \item \textbf{Example:} Euclid proved all truths of 2D geometry from 5 basic axioms.
\end{itemize}

\begin{defbox}[Soundness and Completeness]
\begin{itemize}
  \item \textbf{Sound:} Everything inferred from axioms is true in the model.
  \item \textbf{Complete:} Every true statement about the model can be inferred.
\end{itemize}
Ideally, a logical system is both sound and complete.
\end{defbox}

\paragraph*{Logic in AI}
\begin{itemize}
  \item Designer encodes \textbf{axiomatic truths} $\Delta_0$ about the world in the agent's \textbf{knowledge base} (KB).
  \item Agent uses a \textbf{proof system} (PS) to infer new facts.
  \item Given goals, agent \textbf{acts}, bringing about changes to the world.
  \item Agent \textbf{senses} new facts $\Delta_1$, updates KB, and the cycle continues.
\end{itemize}

\begin{thmbox}[Reductionist Arguments for Logic]
\begin{enumerate}
  \item A complete proof system for first-order logic can simulate all Turing-level computation $\to$ achievement of any complex goal.
  \item If intelligence can be formalized mathematically, and mathematics has logical foundations, then intelligence can be formalized in logic.
\end{enumerate}
\end{thmbox}

\separator

\paragraph*{Modal Logics and BDI}

\begin{defbox}[Modal Logic]
Extends classical logic with \textbf{modalities} that qualify statements:
\begin{itemize}
  \item \textbf{Alethic modalities:} $\Box p$ = necessarily $p$ (true in all possible worlds); $\Diamond p$ = possibly $p$ (true in at least one possible world).
  \item \textbf{Deontic modalities:} $Oq$ = $q$ is obligatory; $Pq$ = $q$ is permitted; $Fq$ = $q$ is forbidden.
\end{itemize}
Deontic logic is relevant to AI ethics (reasoning about what agents \textit{ought} to do).
\end{defbox}

\begin{defbox}[BDI Logic (Belief-Desire-Intention)]
Based on Bratman's model of practical reasoning:
\begin{itemize}
  \item \textbf{Beliefs:} What the agent believes about the world (may not be true).
  \item \textbf{Desires:} Motivational states (what agent would like to accomplish); can be inconsistent.
  \item \textbf{Goals:} Desires being actively pursued; must be consistent.
  \item \textbf{Intentions:} High-level plans to achieve chosen goals.
\end{itemize}

\textbf{Example inference:} $\text{Intend}(q) \to \text{Bel}(\neg q)$ --- if you intend to bring about $q$, you must believe $q$ is not already the case.
\end{defbox}

\bigseparator

% --- Machine Learning ---
\subsubsection{Machine Learning AI}

\begin{defbox}[Machine Learning (Definition)]
``A computer program is said to \textbf{learn from experience E} with respect to some class of tasks T and performance measure P if its performance at tasks in T, as measured by P, \textbf{improves with experience E}.''
\end{defbox}

\begin{center}
\renewcommand{\arraystretch}{1.3}
\begin{tabular}{@{} l p{8cm} @{}}
\toprule
\textbf{Type} & \textbf{Description} \\
\midrule
\textbf{Supervised} & Learns from labeled data (inputs + desired outputs). \\
\textbf{Semi-supervised} & Learns from partially labeled data. \\
\textbf{Unsupervised} & Learns structure from unlabeled data (clustering, etc.). \\
\textbf{Reinforcement} & Learns from feedback (rewards/punishments) in a dynamic environment. \\
\bottomrule
\end{tabular}
\end{center}

\begin{thmbox}[Deep Learning]
\textbf{Deep learning} uses multi-layer neural networks where each layer transforms input into increasingly \textbf{abstract representations}. Responsible for major AI successes:
\begin{itemize}
  \item AlphaGo, AlphaZero (game playing)
  \item Translation, speech recognition
  \item Computer vision, medical imaging
  \item Text generation (GPT-3, etc.)
\end{itemize}
\end{thmbox}

\bigseparator

% --- Limitations ---
\subsubsection{Limitations of Both Paradigms}

\begin{center}
\renewcommand{\arraystretch}{1.3}
\begin{tabular}{@{} p{3.5cm} p{5cm} p{5cm} @{}}
\toprule
\textbf{Issue} & \textbf{Logic-Based AI} & \textbf{Machine Learning} \\
\midrule
Scalability & Too many rules needed; doesn't scale from toy examples to real world. & Needs very large datasets. \\
Brittleness & Small damage $\to$ complete failure. & More robust; degrades gracefully. \\
Sensory tasks & Not designed for low-level perception. & Excels at perception tasks. \\
Symbol grounding & Hand-crafted representations; not grounded in real-world data. & Learns from data; solves grounding. \\
Explainability & Transparent (symbolic). & \textbf{Black box}; opaque reasoning. \\
Hierarchical structure & Natural (recursive rules). & Struggles with hierarchy. \\
Causal reasoning & Can represent causation. & Learns correlation, not causation. \\
Common sense & Theoretically addressed by non-monotonic logic, but computationally hard. & Struggles with common-sense inference. \\
Generalization & Explicit rules generalize. & \textbf{Transfer} problems; poor on novel tasks. \\
Joint reasoning & Not designed for multiple agents. & Hard to integrate with human reasoning. \\
\bottomrule
\end{tabular}
\end{center}

\begin{defbox}[The Symbol Grounding Problem]
In logic-based AI, symbols are \textbf{hand-crafted by designers} rather than grounded in real-world data.

\textbf{Problem:} The semantics (meaning) of symbols relies on the designer's mind, not on direct connection to the world. This prevents ongoing adaptation to unknown environments and is a barrier to full autonomy.

\textbf{ML solution:} Machine learning ``grounds'' representations by extracting them from large datasets.
\end{defbox}

\begin{defbox}[The Black Box Problem]
ML reasoning processes are \textbf{opaque}---even designers can't explain how the system reaches its conclusions.

\textbf{Consequences:}
\begin{itemize}
  \item Cannot provide symbolic/language-like explanations.
  \item Cannot integrate human and machine reasoning.
  \item Has important \textbf{ethical implications} (accountability, trust, bias detection).
\end{itemize}
\end{defbox}

\bigseparator

% --- Non-monotonic Reasoning ---
\subsubsection{Non-Monotonic Reasoning}

\begin{defbox}[Monotonic vs.\ Non-Monotonic Logic]
\textbf{Classical logic is monotonic:}
\[
\text{KB} \cup \Delta_0 \vdash_{\text{PS}} \alpha \implies \text{KB} \cup \Delta_0 \cup \Delta_1 \vdash_{\text{PS}} \alpha
\]
Adding new information never invalidates previous conclusions.

\textbf{Problem:} The real world is too complex to encode all exceptions. We use \textbf{rules of thumb} that can be overridden.

\textbf{Non-monotonic logic:} Conclusions can be \textbf{withdrawn} when new information conflicts with them.
\end{defbox}

\begin{exbox}[Tweety the Penguin]
\begin{itemize}
  \item KB contains: ``Birds typically fly''; ``Tweety is a bird.''
  \item Infer: Tweety flies.
  \item New information: ``Tweety is a penguin.''
  \item Withdraw previous conclusion; infer: Tweety does \textbf{not} fly.
\end{itemize}

The inference $\neg\text{fly}(\text{Tweety})$ is \textbf{preferred} over $\text{fly}(\text{Tweety})$ because penguins are a \textbf{subclass} of birds (specificity principle).
\end{exbox}

\begin{defbox}[Defeasible Rules]
Non-monotonic logics use \textbf{defeasible rules}:
\begin{itemize}
  \item ``Birds \textit{typically}/\textit{usually} fly'' (rather than ``necessarily'').
  \item ``Birds fly \textit{unless there is evidence to the contrary}.''
\end{itemize}
These rules can be overridden by more specific information.
\end{defbox}

\begin{tipbox}[Relevance to AI Ethics]
The problem of exceptions is crucial for ethical AI. You can specify a moral rule (e.g., ``Do not harm humans''), but there will always be exceptions. How do you encode all possible exceptions? Non-monotonic reasoning and argumentation address this.
\end{tipbox}

\bigseparator

% --- Argumentation Theory ---
\subsubsection{Argumentation Theory}

\begin{defbox}[Why Argumentation?]
Non-monotonic reasoning = deciding among \textbf{conflicting inferences} (beliefs, goals, actions).

\textbf{Argumentation} formalizes this as:
\begin{itemize}
  \item More in line with how humans reason (debate, discussion).
  \item Can be formalized to give rational outcomes under limited resources.
  \item Enables \textbf{joint reasoning} among multiple agents (human and AI).
  \item Reasoning becomes \textbf{publicly visible}, not hidden in a black box.
\end{itemize}
\end{defbox}

\begin{defbox}[Arguments and Attacks]
An \textbf{argument} is a pair: $(\text{Premises}, \text{Claim})$.

\textbf{Example:}
\begin{itemize}
  \item $A_1 = (\{\text{penguin}(\text{Tweety}), \text{penguin} \to \neg\text{fly}\}, \neg\text{fly}(\text{Tweety}))$
  \item $A_2 = (\{\text{bird}(\text{Tweety}), \text{bird} \to \text{fly}\}, \text{fly}(\text{Tweety}))$
\end{itemize}

$A_1$ and $A_2$ \textbf{attack} each other (contradictory claims).

If $A_1$ is \textbf{preferred} (specificity: penguins are a subclass of birds), then $A_2$'s attack on $A_1$ is cancelled. $A_1$ wins.
\end{defbox}

\separator

\paragraph*{Dung's Argumentation Framework (1995)}

\begin{defbox}[Argumentation Framework]
An \textbf{argumentation framework} is a pair $\langle \text{Args}, \text{Attacks} \rangle$ where:
\begin{itemize}
  \item $\text{Args}$ = set of arguments.
  \item $\text{Attacks} \subseteq \text{Args} \times \text{Args}$ = binary attack relation.
\end{itemize}

Represented as a \textbf{directed graph}: nodes = arguments, edges = attacks.
\end{defbox}

\begin{defbox}[Labeling Arguments]
Label each argument as:
\begin{itemize}
  \item \textbf{IN} (winning): All attackers are OUT.
  \item \textbf{OUT} (losing): At least one attacker is IN.
  \item \textbf{UNDEC} (undecided): At least one attacker is UNDEC, and no attacker is IN.
\end{itemize}

A \textbf{legal labeling} satisfies these conditions for all arguments.
\end{defbox}

\begin{exbox}[Labeling Example]
Framework: $A_3 \to A_2 \to A_1$ (i.e., $A_3$ attacks $A_2$; $A_2$ attacks $A_1$).

\textbf{Legal labeling:}
\begin{itemize}
  \item $A_3$: IN (no attackers).
  \item $A_2$: OUT (attacker $A_3$ is IN).
  \item $A_1$: IN (all attackers---just $A_2$---are OUT).
\end{itemize}
\end{exbox}

\begin{defbox}[Grounded vs.\ Preferred Extensions]
Given all legal labelings:

\textbf{Grounded extension:}
\begin{itemize}
  \item The unique labeling with \textbf{minimum} IN arguments.
  \item Most skeptical/cautious.
\end{itemize}

\textbf{Preferred extensions:}
\begin{itemize}
  \item Labelings with \textbf{maximal} IN arguments (no superset exists).
  \item May be multiple preferred extensions.
  \item More credulous.
\end{itemize}
\end{defbox}

\begin{exbox}[Mutual Attack: $A_1 \leftrightarrow A_2$]
Three legal labelings:
\begin{enumerate}
  \item $A_1$: IN, $A_2$: OUT --- \textit{preferred}.
  \item $A_1$: OUT, $A_2$: IN --- \textit{preferred}.
  \item $A_1$: UNDEC, $A_2$: UNDEC --- \textit{grounded} (minimum IN = $\emptyset$).
\end{enumerate}

Interpretation: A genuine dilemma. Preferred extensions choose one side; grounded extension remains skeptical.
\end{exbox}

\bigseparator

% --- Argument Games and Dialogue ---
\subsubsection{Argument Games and Dialogue}

\begin{defbox}[Argument Game Proof Theory]
To determine if argument $X$ is in the grounded/preferred extension:
\begin{enumerate}
  \item \textbf{PRO} (proponent) moves $X$.
  \item \textbf{OPP} (opponent) moves an attacking argument.
  \item PRO and OPP alternate, moving attacks according to game rules.
  \item If PRO successfully counters all OPP moves, $X$ is \textbf{justified} (in the extension).
\end{enumerate}

\textbf{Rules differ:}
\begin{itemize}
  \item \textbf{Grounded game:} PRO cannot repeat arguments in a dispute.
  \item \textbf{Preferred game:} OPP cannot repeat arguments in a dispute.
\end{itemize}
\end{defbox}

\separator

\begin{defbox}[From Single-Agent to Distributed Reasoning]
\textbf{Single-agent:} Agent constructs framework from private KB, uses argument game to find winning arguments.

\textbf{Multi-agent (dialogue):} Multiple agents exchange arguments publicly. The framework is built from \textbf{publicly communicated contents}, not private KBs.

\textbf{Public semantics:} Argument $X$ wins iff $X$ is in an extension of the framework built from what has been \textbf{publicly communicated}.
\end{defbox}

\begin{thmbox}[Dialogical Reasoning]
A dialogue that establishes $\alpha$ effectively demonstrates that $\alpha$ can be inferred from the contents of the arguments exchanged during the dialogue.

\textbf{Implications:}
\begin{itemize}
  \item Reasoning is \textbf{transparent} (publicly visible).
  \item Multiple minds bring different knowledge and perspectives.
  \item Enables \textbf{human-AI joint reasoning}.
  \item Important for ensuring \textbf{ethical AI decisions}.
\end{itemize}
\end{thmbox}

\bigseparator

% --- Summary ---
\subsubsection{Summary: Key Takeaways}

\begin{thmbox}[Week 4 Key Points]
\begin{enumerate}
  \item \textbf{Logic-based AI:} Symbolic reasoning; sound/complete proof systems; BDI logics for beliefs, desires, intentions.
  \item \textbf{Machine learning:} Learns from data; deep learning enables abstraction; major successes in perception/generation.
  \item \textbf{Limitations:}
  \begin{itemize}
    \item Logic: Brittleness, symbol grounding, scalability.
    \item ML: Black box, no causation, poor transfer, no joint reasoning.
  \end{itemize}
  \item \textbf{Non-monotonic reasoning:} Conclusions can be withdrawn; uses defeasible rules; addresses real-world exceptions.
  \item \textbf{Argumentation:} Arguments attack each other; preferences determine winners; labelings (IN/OUT/UNDEC).
  \item \textbf{Grounded extension:} Minimum IN (skeptical). \textbf{Preferred extensions:} Maximal IN (credulous).
  \item \textbf{Dialogue:} Enables distributed reasoning; framework built from public communication; transparent, joint AI-human reasoning.
  \item \textbf{Ethical relevance:} Argumentation supports explainable, collaborative reasoning---crucial for ethical AI.
\end{enumerate}
\end{thmbox}

\begin{defbox}[Key Terms to Remember]
\begin{itemize}
  \item \textbf{GOFAI:} Good Old-Fashioned AI (symbolic/logic-based).
  \item \textbf{Soundness/Completeness:} Proof system properties.
  \item \textbf{Modal logic:} Modalities ($\Box$, $\Diamond$, $O$, $P$, $F$).
  \item \textbf{BDI:} Belief-Desire-Intention.
  \item \textbf{Symbol grounding problem:} Symbols not connected to world.
  \item \textbf{Black box problem:} Opaque ML reasoning.
  \item \textbf{Monotonic/Non-monotonic:} Whether new info can retract conclusions.
  \item \textbf{Defeasible rules:} Rules that can be overridden.
  \item \textbf{Argumentation framework:} $\langle \text{Args}, \text{Attacks} \rangle$.
  \item \textbf{IN/OUT/UNDEC:} Argument labels.
  \item \textbf{Grounded extension:} Minimal IN (skeptical).
  \item \textbf{Preferred extension:} Maximal IN (credulous).
  \item \textbf{Dialogical reasoning:} Joint reasoning via public argument exchange.
\end{itemize}
\end{defbox}

